\chapter{Foxglove e interfaccia di visualizzazione}
\label{cap:foxglove}

\begin{notabox}[Contributo extra]
Lo strumento descritto qui non era previsto dal progetto: \`e stato realizzato per
comodit\`a di analisi (in particolare per l'indagine forense dell'incidente del
26/05) e si \`e rivelato utile per tutta la campagna. \`E presentato in fondo come
contributo aggiuntivo.
\end{notabox}

\section{Motivazione}

Gli script matplotlib (cartella \texttt{plot/}) sono ottimi per analisi mirate,
ma non danno una visione \emph{spazio-temporale} del volo. Per ricostruire
visivamente l'incidente serviva un replay 3D sincronizzato con la telemetria.
L'extension PX4 ufficiale di Foxglove si \`e rivelata inaffidabile sui nostri log
(alcuni topic, tra cui \texttt{vehicle\_attitude}, non venivano convertiti,
lasciando vuoto l'albero delle trasformate). La soluzione \`e stata una pipeline
custom che precompila tutto il necessario nel file di output.

\section{Pipeline \texttt{.ulg} $\rightarrow$ \texttt{.mcap}}

Lo script \texttt{foxglove/ulog\_to\_mcap.py} converte il log PX4 in un file
\texttt{.mcap} \textbf{autoconsistente} (nessun plugin richiesto in Foxglove).
Oltre a copiare i 147 topic uORB originali, lo script inietta i dati derivati
necessari al replay 3D:

\begin{itemize}
  \item le trasformate \texttt{/tf} corpo (\texttt{local\_origin}$\rightarrow$
        \texttt{base\_link}) con conversione di coordinate NED$\rightarrow$ENU e
        FRD$\rightarrow$FLU;
  \item le trasformate delle sei eliche, con angolo integrato dagli RPM di
        \texttt{esc\_status} e fattore di scala visivo (per evitare l'aliasing
        \emph{wagon-wheel});
  \item topic interpretativi (\texttt{attitude\_euler} in gradi,
        \texttt{flight\_state} con quota positivo-up e nomi leggibili dei modi di
        volo) e i \texttt{logged\_messages} del firmware.
\end{itemize}

\figurasegnaposto{Schema della pipeline \texttt{.ulg} $\rightarrow$
\texttt{.mcap} $\rightarrow$ Foxglove (topic uORB + transform + derivati).}{Schema
-- da generare}{fig:pipeline-mcap}

\section{Modello 3D e scena}

Il drone \`e rappresentato da un modello \textbf{URDF} puro
(\texttt{f550.urdf}), stilizzato ma fedele alle proporzioni dell'F550 (frame,
mast GPS, sei bracci con motori ed eliche), con i due bracci anteriori colorati
per indicare il fronte. La scelta dell'URDF rispetto a una mesh fotorealistica
privilegia riproducibilit\`a, manutenibilit\`a e leggibilit\`a tecnica. \`E
disponibile anche un \textbf{overlay satellitare} georeferenziato sul punto di
decollo (tile ESRI World Imagery), embeddato direttamente nel \texttt{.mcap}.

\section{Layout dei pannelli}

Il layout committato (\texttt{foxglove/drone f550.json}) organizza 9 pannelli in
due colonne: a sinistra la scena 3D con RPM motori, comandi mixer e input del
pilota; a destra i messaggi di log, le transizioni di stato di volo, la qualit\`a
del fix GPS e l'altitudine. Durante lo scrubbing si vede simultaneamente il drone
muoversi nello spazio e l'evoluzione di tutte le grandezze diagnostiche --- la
modalit\`a con cui \`e stata ricostruita la dinamica dell'incidente del
Capitolo~\ref{cap:diagnostici}.

\figurasegnaposto{Screenshot dell'interfaccia Foxglove realizzata: scena 3D con
modello URDF, overlay satellitare e pannelli di telemetria sincronizzati.}{Screenshot
-- da catturare, chiave}{fig:foxglove-ui}
